\newcommand{\coverpage}{%
  \thispagestyle{empty}
  \setlength{\fboxrule}{0.8pt}
  \setlength{\fboxsep}{0pt}
  \noindent
  \fbox{%
    \begin{minipage}[c][\dimexpr\textheight-2\fboxrule\relax][c]{\dimexpr\textwidth-2\fboxrule\relax}
      \centering
      {\bfseries ĐẠI HỌC QUỐC GIA TP. HỒ CHÍ MINH\par}
      \vspace{0.35cm}
      {\bfseries TRƯỜNG ĐẠI HỌC CÔNG NGHỆ THÔNG TIN\par}
      \vspace{0.35cm}
      {\bfseries KHOA HỆ THỐNG THÔNG TIN\par}

      \vspace{2.9cm}
      {\bfseries NGÔ TIẾN SỸ -- 23521367\par}
      \vspace{0.2cm}
      {\bfseries NGUYỄN VĂN NAM -- 23520982\par}

      \vspace{2.7cm}
      {\bfseries BÁO CÁO ĐỒ ÁN\par}
      \vspace{0.55cm}
      {\bfseries
      \begin{minipage}{0.86\textwidth}
        \centering
        XÂY DỰNG KIẾN TRÚC KHO HỒ DỮ LIỆU THỜI GIAN THỰC CHO HỆ THỐNG GIAO THÔNG THÔNG MINH TRÊN CỤM KUBERNETES
      \end{minipage}\par}
      \vspace{0.5cm}
      {\bfseries\color{red}
      \begin{minipage}{0.86\textwidth}
        \centering
        Building a Real-Time Data Lakehouse Architecture for Intelligent Transportation Systems on a Kubernetes Cluster
      \end{minipage}\par}

      \vspace{2.2cm}
      {\bfseries CHUYÊN NGÀNH HỆ THỐNG THÔNG TIN\par}

      \vspace{2.1cm}
      {\bfseries GIẢNG VIÊN HƯỚNG DẪN\par}
      \vspace{0.35cm}
      {\bfseries THS. NGUYỄN HỒ DUY TRÍ\par}

      \vfill
      {\bfseries TP. HỒ CHÍ MINH, 2026\par}
      \vspace{0.6cm}
    \end{minipage}%
  }
  \par
  \clearpage
}

\newcommand{\acknowledgements}{%
  \thispagestyle{empty}
  \begin{center}
    {\bfseries LỜI CẢM ƠN}
  \end{center}
  \vspace{1cm}
  Nhóm thực hiện xin trân trọng cảm ơn ThS. Nguyễn Hồ Duy Trí đã hướng dẫn, góp ý và định hướng trong quá trình xây dựng đề tài. Những trao đổi về kiến trúc dữ liệu, xử lý luồng, triển khai trên Kubernetes và cách diễn giải kết quả thực nghiệm giúp nhóm hoàn thiện báo cáo theo hướng rõ ràng và có khả năng kiểm chứng hơn.

  Nhóm cũng xin cảm ơn Trường Đại học Công nghệ Thông tin, Đại học Quốc gia TP. Hồ Chí Minh và Khoa Hệ thống Thông tin đã tạo môi trường học tập để nhóm có thể tiếp cận các kiến thức về cơ sở dữ liệu phân tán, dữ liệu lớn, hạ tầng Kubernetes, GitOps và quan sát hệ thống. Các kiến thức này là nền tảng để nhóm thiết kế, triển khai và đánh giá hiện vật Continux trong phạm vi đồ án.

  Do giới hạn thời gian và hạ tầng thực nghiệm, báo cáo khó tránh khỏi thiếu sót. Nhóm mong nhận được góp ý từ giảng viên và các bạn để tiếp tục cải thiện hệ thống trong các phiên bản sau.

  \vspace{1.5cm}
  \begin{flushright}
  \begin{tabular}{c}
  TP. Hồ Chí Minh, 2026\\
  Nhóm thực hiện
  \end{tabular}
  \end{flushright}
  \clearpage
}

\newcommand{\abbreviationlist}{%
  \thispagestyle{empty}
  \begin{center}
    {\bfseries DANH MỤC TỪ VIẾT TẮT}
  \end{center}
  \vspace{0.8cm}
  \begin{tabularx}{\textwidth}{@{}p{0.18\textwidth}Y@{}}
  \toprule
  \textbf{Từ viết tắt} & \textbf{Diễn giải}\\
  \midrule
  API & Application Programming Interface -- giao diện lập trình ứng dụng\\
  CSV & Comma-Separated Values -- định dạng dữ liệu phân tách bằng dấu phẩy\\
  DSR & Design Science Research -- phương pháp nghiên cứu khoa học thiết kế\\
  GitOps & Mô hình vận hành dùng Git làm nguồn trạng thái mong muốn\\
  HA & High Availability -- tính sẵn sàng cao\\
  ITS & Intelligent Transportation System -- hệ thống giao thông thông minh\\
  JSONL & JSON Lines -- định dạng mỗi dòng là một JSON độc lập\\
  K3s & Bản phân phối Kubernetes gọn nhẹ của SUSE\\
  MV & Materialized View -- khung nhìn vật lý\\
  PVC & PersistentVolumeClaim -- yêu cầu cấp phát lưu trữ bền vững trong Kubernetes\\
  S3 & Simple Storage Service -- giao diện lưu trữ đối tượng tương thích Amazon S3\\
  SDG & Sustainable Development Goal -- Mục tiêu Phát triển Bền vững\\
  SQL & Structured Query Language -- ngôn ngữ truy vấn có cấu trúc\\
  TLC & Taxi and Limousine Commission -- Ủy ban Taxi và Limousine Thành phố New York\\
  \bottomrule
  \end{tabularx}
  \clearpage
}

\newcommand{\reportabstract}{%
  \clearpage
  \hypersetup{pageanchor=true}
  \pagenumbering{arabic}
  \setcounter{page}{1}
  \pagestyle{fancy}
  \phantomsection
  \addcontentsline{toc}{section}{Tóm tắt báo cáo}
  \begin{center}
    {\bfseries TÓM TẮT BÁO CÁO}
  \end{center}
  \vspace{0.5cm}
  Nghiên cứu này xây dựng và đánh giá \textbf{Continux}, một kiến trúc kho hồ dữ liệu thời gian thực (\emph{real-time Data Lakehouse}) cho dữ liệu giao thông thông minh trên cụm K3s ba máy chủ. Hệ thống chuyển đổi dữ liệu NYC TLC Yellow Taxi từ Parquet sang JSONL, phát lại từng bản ghi thành sự kiện bằng Vector, chuyển luồng qua Redpanda, xử lý bằng SQL luồng trong RisingWave và ghi kết quả Apache Iceberg trên MinIO. Trạng thái triển khai được quản lý bằng Argo CD; chỉ số vận hành được lưu bằng VictoriaMetrics và trực quan hóa bằng Grafana. Hiện vật kỹ thuật được đánh giá qua bốn lượt chạy độc lập với giới hạn phát cấu hình lần lượt là \texttt{2}, \texttt{1.000}, \texttt{5.000} và \texttt{10.000} sự kiện mỗi giây. Ở cả bốn lượt, đầu ra Iceberg được tạo, Vector trở về trạng thái dừng sau phát lại và vòng lặp truy vấn không ghi nhận lỗi trong lúc hoán đổi khung nhìn vật lý công khai. Lượt \texttt{benchmark-high} ghi nhận \texttt{426.960} chuyến trước chuyển giao, \texttt{423.239} chuyến sau khi áp dụng cách lọc mới và thời lượng hoán đổi \texttt{0,215989 s}. Kết quả chứng minh khả năng cập nhật cách tính truy vấn công khai mà không ghi nhận gián đoạn trong cửa sổ quan sát, đồng thời cho thấy quy trình thực nghiệm có thể lặp lại trên cụm tài nguyên giới hạn.

  \vspace{0.5cm}
  \noindent\textbf{Từ khóa:} kho hồ dữ liệu, xử lý luồng, khung nhìn vật lý, chuyển giao Xanh lam/Xanh lục, RisingWave, Apache Iceberg, GitOps, K3s, hệ thống giao thông thông minh.
  \clearpage
}

\pagenumbering{gobble}
\hypersetup{pageanchor=false}
\newgeometry{top=1cm,bottom=1cm,left=1cm,right=1cm}
\coverpage
\coverpage
\restoregeometry
\acknowledgements

\thispagestyle{empty}
\tableofcontents
\clearpage

\thispagestyle{empty}
\listoffigures
\clearpage

\thispagestyle{empty}
\listoftables
\clearpage

\abbreviationlist
\reportabstract
